\chapter{Ray Tracing ISA Extensions for Borg}
\label{chap:isa}

\section{Motivation}

Standard rasterisation pipelines are ill-suited for rays:
a ray-scene intersection test requires traversing a hierarchical
acceleration structure (BVH) whose depth and branching factor are
data-dependent.
NVIDIA addressed this with dedicated RT Cores in Turing; we address it
via ISA extensions to \borg{}'s existing scalar shader pipeline.

\section{BVH Fundamentals}

\todo{BVH construction (SAH), traversal algorithm, stack management.}

\section{New Instructions}

\begin{table}[h]
  \centering
  \begin{tabular}{lll}
    \toprule
    Mnemonic & Operands & Description \\
    \midrule
    \texttt{RAY.INIT}   & rd, rs1, rs2 & Initialise ray (origin + direction) \\
    \texttt{BVH.PUSH}   & rs1          & Push BVH node onto traversal stack  \\
    \texttt{BVH.TEST}   & rd, rs1      & AABB intersection test              \\
    \texttt{TRI.ISECT}  & rd, rs1      & Möller–Trumbore triangle intersection \\
    \texttt{RAY.RESULT} & rd           & Read closest hit record             \\
    \bottomrule
  \end{tabular}
  \caption{Ray tracing ISA extensions.}
  \label{tab:isa_rt}
\end{table}

\section{Hardware Implementation}

\subsection{AABB Intersection Unit}

\todo{Parametric slab method. Floating-point pipeline. Latency.}

\subsection{Möller–Trumbore Unit}

\todo{Triangle intersection. Edge vectors, barycentric coordinates.}

\subsection{Traversal Stack}

\todo{Dedicated BRAM-backed stack. Depth limit. Stack overflow handling.}

\section{Integration with the Shader Pipeline}

\todo{How new instructions slot into the existing scalar execution unit.
      Stall/forwarding behaviour.}

\section{Summary}

\todo{Area overhead estimate. Expected throughput.}
